\documentclass[a4paper,12pt]{article}\usepackage{tikz} \usepackage[utf8]{inputenc} \usepackage[romanian]{babel} \usepackage[T1]{fontenc} \usepackage{amsmath} \usepackage{geometry} \geometry{margin=2.5cm} \title{Inteligență Artificială\\Tutoriatul 1} \author{} \date{} \begin{document} \maketitle \section{Introducere} Inteligența Artificială (IA) studiază construirea unor sisteme capabile să execute sarcini care necesită în mod normal inteligență umană. \medskip Principalele direcții studiate în IA sunt: \begin{itemize} \item reprezentarea cunoștințelor; \item raționamentul automat; \item căutarea în spațiul stărilor; \item învățarea automată; \item procesarea limbajului natural; \item vederea artificială. \end{itemize} \section{Agentul inteligent} Un agent este o entitate care percepe mediul și acționează asupra lui. \bigskip Schema generală: \[ Mediu \rightarrow Senzori \rightarrow Agent \rightarrow Actuatori \rightarrow Mediu \] Exemple de agenți: \begin{itemize} \item robot aspirator; \item mașină autonomă; \item sistem de recomandare; \item ChatGPT. \end{itemize} \section{Proprietățile mediului} \subsection{Complet observabil vs. parțial observabil} \textbf{Complet observabil} \begin{itemize} \item agentul cunoaște întreaga stare a mediului; \item exemplu: șah. \end{itemize} \textbf{Parțial observabil} \begin{itemize} \item informația disponibilă este incompletă; \item exemplu: poker. \end{itemize} \subsection{Determinist vs. stochastic} \textbf{Determinist} Aceeași acțiune produce întotdeauna același rezultat. \textbf{Stochastic} Există elemente de aleator. \subsection{Static vs. dinamic} \textbf{Static} Mediul nu se modifică în timpul luării deciziei. \textbf{Dinamic} Mediul se poate schimba în timp. \subsection{Discret vs. continuu} \textbf{Discret} Număr finit de stări și acțiuni. \textbf{Continuu} Valorile pot varia continuu. \section{Probleme de căutare} O problemă de căutare este definită prin: \begin{enumerate} \item starea inițială; \item starea scop; \item mulțimea acțiunilor; \item funcția succesor; \item costul acțiunilor. \end{enumerate} \subsection{Exemplu} Dorim să găsim un drum de la nodul A la nodul G.
\begin{center} \begin{tikzpicture}[ every node/.style={circle,draw,minimum size=8mm} ] \node (A) at (0,0) {A}; \node (B) at (-2,-2) {B}; \node (C) at ( 2,-2) {C}; \node (D) at (-3,-4) {D}; \node (E) at (-1,-4) {E}; \node (F) at ( 2,-4) {F}; \node (G) at (-1,-6) {G}; \draw (A) -- (B); \draw (A) -- (C); \draw (B) -- (D); \draw (B) -- (E); \draw (C) -- (F); \draw (E) -- (G); \end{tikzpicture} \end{center}
 \section{Graf de stări vs. arbore de căutare} \subsection{Graf de stări} O stare apare o singură dată. \subsection{Arbore de căutare} Aceeași stare poate apărea de mai multe ori deoarece putem ajunge la ea pe drumuri diferite. \bigskip Această diferență este foarte importantă la examen. \section{Frontiera} Frontiera reprezintă mulțimea nodurilor generate dar neexplorate încă. \bigskip Diferența dintre algoritmi constă în modul în care aleg următorul nod din frontieră. \section{Breadth First Search (BFS)} \subsection{Idee} Explorează nodurile pe niveluri. \bigskip Structură utilizată: \[ \textbf{Queue (FIFO)} \] Primul introdus este primul extras. \subsection{Aplicare pe exemplu} Ordinea de explorare: \begin{verbatim} A B C D E F G \end{verbatim} Drumul găsit: \begin{verbatim} A -> B -> E -> G \end{verbatim} \subsection{Proprietăți} \begin{itemize} \item complet; \item optimal pentru costuri egale; \item consum mare de memorie. \end{itemize} \section{Depth First Search (DFS)} \subsection{Idee} Explorează cât mai adânc înainte de a reveni. \bigskip Structură utilizată: \[ \textbf{Stack (LIFO)} \] Ultimul introdus este primul extras. \subsection{Aplicare pe exemplu} Ordinea de explorare: \begin{verbatim} A B D E G \end{verbatim} \subsection{Proprietăți} \begin{itemize} \item memorie redusă; \item nu este optimal; \item poate intra pe ramuri foarte adânci; \item nu este complet în spații infinite. \end{itemize} \section{Depth Limited Search (DLS)} \subsection{Idee} DFS cu o limită maximă de adâncime. \bigskip Presupunem limita: \[ l = 2 \] și drumul soluție: \begin{verbatim} A -> B -> E -> G \end{verbatim} Adâncimea soluției este: \[ 3 \] Algoritmul nu va ajunge la G deoarece nu poate coborî sub adâncimea 2. \subsection{Concluzie} DLS poate rata soluția dacă aceasta se află mai jos decât limita impusă. \section{Iterative Deepening Search (IDS)} \subsection{Idee} Execută DLS cu limite crescătoare: \begin{verbatim} 0 1 2 3 ... \end{verbatim} Pentru exemplul nostru: \begin{itemize} \item limita 0: nu găsește soluția; \item limita 1: nu găsește soluția; \item limita 2: nu găsește soluția; \item limita 3: găsește soluția. \end{itemize} \subsection{Avantaje} \begin{itemize} \item complet; \item optimal pentru costuri egale; \item consum redus de memorie. \end{itemize} \section{Uniform Cost Search (UCS)} Considerăm acum costurile: \begin{verbatim} A -> B = 1 A -> C = 5 B -> D = 2 B -> E = 1 E -> G = 1 C -> F = 1 \end{verbatim} \subsection{Idee} Algoritmul alege întotdeauna nodul cu costul total minim. \bigskip Structură utilizată: \[ \textbf{Priority Queue} \] \subsection{Calcul} Costul drumului: \begin{verbatim} A -> B -> E -> G \end{verbatim} este: \[ 1 + 1 + 1 = 3 \] Acesta va fi drumul ales. \subsection{Observație} BFS minimizează numărul de pași. UCS minimizează costul total. \section{Comparație între algoritmi} \begin{center} \begin{tabular}{|c|c|} \hline Algoritm & Structură utilizată \\ \hline BFS & Queue (FIFO) \\ DFS & Stack (LIFO) \\ DLS & Stack + limită \\ IDS & DLS repetat \\ UCS & Priority Queue \\ \hline \end{tabular} \end{center} \section{Întrebări de verificare} \begin{enumerate} \item Ce este un agent inteligent? \item Care este diferența dintre mediu determinist și stochastic? \item Care este diferența dintre graf de stări și arbore de căutare? \item Ce reprezintă frontiera? \item Ce structură de date utilizează BFS? \item Ce structură de date utilizează DFS? \item De ce poate DLS rata soluția? \item Care este avantajul principal al IDS? \item Care este diferența dintre BFS și UCS? \item În ce situații UCS este preferat în locul BFS? \end{enumerate} \end{document}